% !TeX program = pdflatex
% =============================================================================
% Arbeitsblatt: Bewegungsgleichungen (Teil 2) (nicht bewertet, Übungsmaterial
% für den Unterricht)
% Kantonsschule KFR – Physik, 4. Klasse, 1. HS, Mechanik
%
% Fortsetzung von bewegungsgleichungen-arbeitsblatt.tex (Teil 1): Teil 1
% deckt die Herleitung aller drei Bewegungsgleichungen, die Formelwahl
% (Checkliste "Wie wähle ich die passende Bewegungsgleichung?"), die
% Vorzeichenkonvention bei Bremsbewegungen (E-Trottinett-Beispiel) und die
% Treffpunktprobleme mit Beschleunigung ab (dieser letzte Block stand
% ursprünglich am Ende des ungeteilten Dokuments, wurde aber zu Teil 1
% verschoben, weil dort auch das dafür nötige Vorwissen -- Gleichsetzen,
% Formelwahl -- liegt). Teil 2 baut direkt darauf auf und vertieft die
% Anwendung: Anhalteweg (Reaktionsweg/Bremsweg), Einfluss von Reaktionszeit
% und Untergrund, Faustregeln aus der Fahrschule sowie quadratische
% Gleichungen (Mitternachtsformel).
%
% Eigenständige Kopie derselben Farbdefinitionen und Boxdefinitionen wie
% Teil 1 (Templates sind selbstständig/nicht DRY, siehe CLAUDE.md).
%
% Stand: Theorie und Aufgaben aus dem vormals ungeteilten Arbeitsblatt
% übernommen (unverändert bis auf zwei Verweise, die zuvor auf Inhalte
% "weiter oben" im selben Dokument zeigten und jetzt auf "Teil 1"
% umformuliert sind, weil der Inhalt zur Parabel und zur
% Formelwahl-Checkliste, auf den sie zeigen, dort liegt). Lernziele und
% Einleitung neu für Teil 2 geschrieben, Nummerierung LZ1..LZ7 folgt der
% Dokumentreihenfolge dieses Teils (eigenständig, nicht identisch mit der
% Nummerierung in Teil 1 oder in lernziele.md).
% =============================================================================
\documentclass[addpoints,11pt,a4paper]{exam}

\usepackage[T1]{fontenc}
\usepackage[utf8]{inputenc}
\usepackage[ngerman]{babel}
\usepackage{lmodern}
\usepackage[a4paper,margin=2.2cm,headheight=14pt]{geometry}
\usepackage{amsmath}
\usepackage{siunitx}
% Hausstil: zusammengesetzte Einheiten immer als echter Bruch (\frac), nicht
% als Schrägstrich inline oder \tfrac. Gilt global für jedes \si/\SI.
\sisetup{per-mode=fraction,fraction-command=\frac}
\usepackage{array,booktabs,tabularx,multirow}
\usepackage{enumitem}
\usepackage{graphicx}
\usepackage{xcolor}
\usepackage{colortbl}
\usepackage{tikz}
\usepackage{physics}
\usepackage{circuitikz}
\usepackage[most]{tcolorbox}
\usepackage{lastpage}
\usepackage{needspace}

\usetikzlibrary{arrows.meta,calc,angles,quotes}

\setlength{\parindent}{0pt}
\setlength{\parskip}{0.45em}

% -----------------------------
% Farben (Hausfarbschema Physik KFR)
% -----------------------------
\definecolor{titleblue}{HTML}{183A6B}
\definecolor{accentblue}{HTML}{2E6F9E}
\definecolor{lightblue}{HTML}{EAF4FB}
\definecolor{verylightblue}{HTML}{F7FBFE}
\definecolor{lightgray}{HTML}{F4F6F8}
\definecolor{darkgray}{HTML}{4A4A4A}
% Für Diagramme mit zwei verschiedenen Funktionen/Linien: titleblue (dunkles
% Blau) + accentorange (Orange) sind weit genug auseinander (auch in
% Graustufen/für Farbenblinde unterscheidbar), im Gegensatz zu
% titleblue+accentblue, die beide Blautöne sind.
\definecolor{accentorange}{HTML}{D9720C}
% Für Diagramme mit DREI Linien: dritte Farbe nötig, die sich sowohl von
% titleblue als auch von accentorange klar unterscheidet, auch in
% Graustufen/für Farbenblinde. Violett liegt weit genug von Blau und Orange
% entfernt, um als dritte Linie eindeutig unterscheidbar zu bleiben.
\definecolor{accentpurple}{HTML}{7B2D8E}
% Für die Anhalteweg-Diagramme: Reaktionsweg immer titleblue (Blau),
% Hindernis immer accentorange (Orange, bereits definiert). Für den
% Bremsweg (Rot) braucht es eine eigene, klar von Blau/Orange
% unterscheidbare dritte Farbe -- accentpurple ist hier ungeeignet
% (steht schon für "dritte Linie" in anderen Diagrammen), deshalb ein
% eigenes gedecktes Rot.
\definecolor{brakered}{HTML}{C0392B}
% Kein eigenes Grün mehr für die Theoriebox: siehe Kommentar am
% Dateianfang, das Boxen-Layout wurde auf das ungerahmte Layout von
% beschleunigung-arbeitsblatt.tex umgestellt.
% Gruppenarbeitbox: eigene Farbe (Petrol/Teal) bleibt bewusst erhalten
% (auch wenn taskbox/theoriebox jetzt ungerahmt sind), damit Partnerarbeit
% weiterhin auf einen Blick erkennbar ist. Zusätzlich als Linienfarbe in
% der Glatteis-Parabel weiter unten im Dokument wiederverwendet.
\definecolor{groupteal}{HTML}{1B7F72}
\definecolor{lightgroupteal}{HTML}{E8F5F3}

% -----------------------------
% Spaltentypen
% -----------------------------
\newcolumntype{Y}{>{\centering\arraybackslash}X}
\newcolumntype{P}[1]{>{\raggedright\arraybackslash}p{#1}}

% -----------------------------
% Boxen
% -----------------------------
\tcbset{before skip=10pt, after skip=10pt}
\newtcolorbox{titlebox}{
  enhanced,
  colback=titleblue,
  colframe=titleblue,
  boxrule=0pt,
  arc=3mm,
  left=3mm,right=3mm,top=3mm,bottom=3mm
}

% infobox/theoriebox/beispielbox: ungerahmter Fliesstext statt farbiger Box
% (Layout wie beschleunigung-arbeitsblatt.tex / institutionen/UZH-Praktikum-I),
% Titel wird -- wo vorhanden -- manuell/über das optionale Argument fett
% gesetzt.
% Regel bleibt unverändert: hervorgehobene/definierte Begriffe INNERHALB
% einer theoriebox immer mit \textbf{\emph{...}} setzen (fett UND kursiv),
% nicht mit blossem \emph{...} -- damit Fachbegriffe beim Durchlesen der
% Theorie sofort auffallen. Ausserhalb von theoriebox (beispielbox,
% Aufgabentext, ...) bleibt \emph{...} allein korrekt.
\newenvironment{infobox}{%
  \par\addvspace{10pt}\noindent
}{%
  \par\addvspace{10pt}
}

\newenvironment{theoriebox}[1][Theorie]{%
  \par\addvspace{10pt}\noindent\textbf{#1}\par\smallskip\noindent
}{%
  \par\addvspace{10pt}
}

% taskbox/groupbox bleiben technisch tcolorboxen (taskbox davon unsichtbar
% gemacht: keine Farbe/kein Rahmen/kein Padding), weil sie innerhalb der
% exam-Klasse questions-Liste vor dem ersten \question stehen können
% (Bild/Fliesstext) -- eine reine \newenvironment löst dort "Something's
% wrong--perhaps a missing \item", sobald z.B. ein \begin{center} vor dem
% ersten \item der Aussenliste steht.
% exam.cls rückt die questions-Liste um die Breite von "10." (plus
% \labelsep) ein (Platz für die Nummerierung); ohne Korrektur würde der
% Inhalt von taskbox/groupbox (Diagramme, Fliesstext) deshalb sichtbar
% weiter eingerückt erscheinen als der Inhalt von infobox/theoriebox/
% beispielbox (die ausserhalb von questions stehen und deshalb bereits die
% volle Randbreite nutzen) -- gleicher Randabstand im ganzen Dokument, auch
% bei den Übungen.
\newlength{\aufgabenindent}
\settowidth{\aufgabenindent}{10.\hskip\labelsep}
\newtcolorbox{taskbox}[1][]{
  enhanced, breakable,
  colback=white, colframe=white, boxrule=0pt,
  left=0mm,right=0mm,top=0mm,bottom=0mm,
  before skip=10pt, after skip=10pt,
  enlarge left by=-\aufgabenindent,
  width=\textwidth,
  before upper={%
    \def\tmpboxtitle{#1}%
    \ifx\tmpboxtitle\empty\else\noindent\textbf{#1}\par\smallskip\fi
  }
}

\newenvironment{beispielbox}[1][Beispiel]{%
  \par\addvspace{10pt}\noindent\textbf{#1}\par\smallskip\noindent
}{%
  \par\addvspace{10pt}
}

% Regel für groupbox: JEDE Aufgabe, bei der Schülerinnen/Schüler zu zweit
% oder in einer Gruppe zusammenarbeiten (Partnerarbeit, Gruppenarbeit,
% Diskussion zu zweit, ...), bekommt diese Box statt der (jetzt ungerahmten)
% taskbox. groupbox bleibt bewusst farbig (Petrol/Teal, siehe Kommentar bei
% den Farbdefinitionen oben), damit "hier braucht es eine Partnerin/einen
% Partner" trotz des sonst ungerahmten Layouts auf einen Blick erkennbar
% bleibt. Strukturell wie taskbox innerhalb von questions (gleiche
% Einrückungskorrektur).
\newtcolorbox{groupbox}[1][Gruppenarbeit]{
  enhanced,
  breakable,
  colback=lightgroupteal,
  colframe=groupteal,
  title={#1},
  fonttitle=\bfseries,
  boxrule=0.7pt,
  arc=2mm,
  enlarge left by=-\aufgabenindent,
  width=\textwidth,
  left=5mm,right=5mm,top=2mm,bottom=2mm
}

% hinweisbox: eigene Farbe (Orange, wiederverwendet accentorange statt einer
% neuen Farbe), für wichtige Klarstellungen, die über eine normale
% Bemerkung im Fliesstext (wie "\textbf{Wichtig:}" innerhalb einer
% theoriebox) hinausgehen, etwa eine explizite Warnung vor einer häufigen
% Verwechslung. Bleibt trotz des sonst ungerahmten Layouts farbig, damit
% sie weiterhin auffällt.
\definecolor{lightorange}{HTML}{FCEEE0}
\newtcolorbox{hinweisbox}[1][Hinweis]{
  enhanced,
  breakable,
  colback=lightorange,
  colframe=accentorange,
  title={#1},
  fonttitle=\bfseries,
  boxrule=0.7pt,
  arc=2mm,
  left=5mm,right=5mm,top=2mm,bottom=2mm
}

% formelbox: eigene Farbe (Violett, wiederverwendet accentpurple), um die
% drei zentralen Bewegungsgleichungen dieser Einheit visuell klar von
% normaler Theorie abzuheben, damit sie beim Nachschlagen sofort ins Auge
% springen. Steht wie infobox/theoriebox/beispielbox ausserhalb von
% questions, KEINE Einrückungskorrektur nötig (siehe Kommentar bei infobox
% oben).
\definecolor{lightpurple}{HTML}{F3EAF7}
\newtcolorbox{formelbox}[1][Bewegungsgleichung]{
  enhanced,
  breakable,
  colback=lightpurple,
  colframe=accentpurple,
  title={#1},
  fonttitle=\bfseries,
  boxrule=0.9pt,
  arc=2mm,
  left=5mm,right=5mm,top=2mm,bottom=2mm
}

% anleitungbox: eigene Farbe (Gold/Bernstein), für Schritt-für-Schritt-
% Anleitungen/Checklisten (wie die Formelwahl weiter unten), die sich von
% einer normalen theoriebox (ungerahmt), einer hinweisbox (Orange, für
% Klarstellungen) und einer formelbox (Violett, für die Kernformeln)
% farblich unterscheiden sollen.
\definecolor{strategygold}{HTML}{9C7A17}
\definecolor{lightstrategygold}{HTML}{FBF3DE}
\newtcolorbox{anleitungbox}[1][Anleitung]{
  enhanced,
  breakable,
  colback=lightstrategygold,
  colframe=strategygold,
  title={#1},
  fonttitle=\bfseries,
  boxrule=0.7pt,
  arc=2mm,
  left=5mm,right=5mm,top=2mm,bottom=2mm
}

\newtcolorbox{answerbox}{
  breakable,
  colback=white,
  colframe=gray!55,
  boxrule=0.5pt,
  arc=1.5mm,
  left=2mm,right=2mm,top=2mm,bottom=2mm
}

% --- Lösungen ein-/ausblenden -----------------------------------------------
% Schülerexemplar (Standard): \noprintanswers
% Lehrer-/Lösungsexemplar:    \printanswers
%\noprintanswers
\printanswers

% --- Punkte bleiben intern, werden aber nicht gedruckt ----------------------
% Arbeitsblätter sind nicht benotet. \question[n]/\part[n] bekommen trotzdem
% eine Punktzahl n, rein als interner Anhaltspunkt, der 1:1 an
% \Antwortraum{n} weitergereicht wird, um die Antwortraumgrösse zu
% bestimmen. \pointformat{} (Aufruf, nicht \renewcommand!) unterdrückt die
% von der exam-Klasse sonst automatisch gedruckte Punktanzeige ("(n P.)")
% vollständig, ohne die interne Punktezählung (\numpoints) zu beeinflussen.
\pointformat{}

% --- Lösungsbox auf Deutsch, farblich abgesetzt -----------------------------
\renewcommand{\solutiontitle}{\noindent\color{titleblue}\textbf{Lösung:}\enspace}

% --- Kopf-/Fusszeile: Thema/Titel oben, Seitenzahl unten --------------------
\pagestyle{headandfoot}
\cfoot{\textcolor{darkgray}{Seite \thepage\ von \pageref{LastPage}}}

% --- Kopfzeile: Klasse / Datum / Thema / Titel ------------------------------
% Einheitliches Layout für Arbeitsblatt, Prüfung und Praktikum. Ohne Tabelle:
% Klasse/Datum/Thema/Lehrperson stehen als einfache Textzeile unter dem
% Titel-Banner. Fünftes Argument (Untertitel): wie in
% beschleunigung-arbeitsblatt.tex eine rein optische Ergänzung im
% Titel-Banner, die vier Kernangaben Klasse/Datum/Thema/Titel bleiben in
% derselben Reihenfolge/Position wie im Hausstil-Standard.
\newcommand{\Kopfzeile}[5]{%
  \begin{titlebox}
    {\color{white}\sffamily\bfseries\Large #4}\\[3pt]
    {\color{white}\sffamily\normalsize #5}
  \end{titlebox}
  \vspace{0.5em}
  \noindent\textbf{Klasse:}~#1 \hfill \textbf{Datum:}~#2 \hfill \textbf{Thema:}~#3
    \hfill \textbf{Lehrperson:}~Loris Keller
  \vspace{0.8em}
  \lhead{\textcolor{darkgray}{#3}}
  \rhead{\textcolor{darkgray}{#4}}
}

% --- Antwortraum: ca. 2.5 cm Platz pro Punkt, als umrandetes Feld -----------
\newlength{\antwortraumlen}
\newcommand{\Antwortraum}[1]{%
  \pgfmathsetlength{\antwortraumlen}{#1*2.5cm}%
  \begin{answerbox}
    \rule{0pt}{\antwortraumlen}%
  \end{answerbox}%
}


% Werte für Kopfzeile zentral setzen
\newcommand{\Klasse}{4a}
\newcommand{\Datum}{TT.MM.JJJJ}
\newcommand{\Thema}{Bewegungsgleichungen (Teil 2)}
\newcommand{\ArbeitsblattTitel}{Arbeitsblatt: Bewegungsgleichungen (Teil 2)}
\newcommand{\ArbeitsblattUntertitel}{Anhalteweg, Faustregeln und quadratische Gleichungen}

\begin{document}

\Kopfzeile{\Klasse}{\Datum}{\Thema}{\ArbeitsblattTitel}{\ArbeitsblattUntertitel}

% --- Lernziele: neu für Teil 2 formuliert, Reihenfolge = Dokumentreihenfolge
% dieses Teils (Anhalteweg -> s-v-Diagramm -> Reaktionszeit -> Untergrund ->
% Faustregeln -> kritische Geschwindigkeit -> Mitternachtsformel). Das
% Vorzeichen-Lernziel (Bremsbewegung, E-Trottinett) und das
% Treffpunktprobleme-Lernziel wurden nach Teil 1 verschoben, zusammen mit
% den zugehörigen Beispielen/Übungen; deshalb hier nicht mehr aufgeführt. ---
\begin{infobox}
\textbf{Lernziele}\\
Am Ende dieser Einheit kann ich \dots
\begin{itemize}[leftmargin=1.2cm,itemsep=0.25em]
  \item den Anhalteweg als Summe aus Reaktionsweg und Bremsweg berechnen und erklären, welche Rolle die Reaktionszeit dabei spielt. (LZ1)
  \item am $s$-$v$-Diagramm erklären, weshalb der Reaktionsweg linear und der Bremsweg quadratisch mit $v_0$ wächst, und daraus die höchste sichere Geschwindigkeit vor einem Hindernis ablesen. (LZ2)
  \item berechnen, wie stark eine verlängerte Reaktionszeit (zum Beispiel unter Alkoholeinfluss) den Anhalteweg vergrössert. (LZ3)
  \item den Bremsweg auf verschiedenen Untergründen vergleichen und erklären, weshalb Aquaplaning und Glatteis besonders gefährlich sind. (LZ4)
  \item bekannte Faustregeln aus der Fahrschule (2-Sekunden-Regel, Halber-Tacho-Regel, Bremswegfaustformel) mit dem exakt berechneten Anhalteweg vergleichen und ihre Grenzen benennen. (LZ5)
  \item aus einer gegebenen Bremsstrecke die kritische Anfangsgeschwindigkeit mit der Mitternachtsformel berechnen und die unphysikalische Lösung verwerfen. (LZ6)
  \item eine quadratische Bewegungsgleichung mit der Mitternachtsformel nach der Zeit auflösen und unter mehreren mathematischen Lösungen die physikalisch sinnvolle auswählen. (LZ7)
\end{itemize}
\end{infobox}

% --- Einleitung: Anhalteweg als konkrete Alltagsgefahr, direkt an Teil 1
% anschliessend, noch ohne die Anhalteweg-Formel selbst zu nennen (die folgt
% erst in der ersten Theoriebox). ---------------------------------------
\begin{beispielbox}[Einleitung: Warum der Anhalteweg über Leben und Tod entscheidet]
Stell dir vor, du fährst Velo oder sitzt im Auto, als plötzlich ein Ball
auf die Strasse rollt und ihm ein Kind hinterherläuft. Ab dem Moment, in
dem du das Kind siehst, bis das Fahrzeug tatsächlich steht, vergeht
wertvolle Zeit und Strecke, selbst bei einer sofortigen Vollbremsung.

Genau diese Strecke, der Anhalteweg, entscheidet in einer solchen
Situation, ob ein Unfall verhindert werden kann. Er hängt nicht nur davon
ab, wie stark gebremst wird, sondern auch davon, wie schnell die fahrende
Person überhaupt reagiert und wie griffig die Strasse ist. In dieser
Einheit berechnest du mit den Bewegungsgleichungen aus Teil 1 exakt, wie
sich diese Faktoren auf den Anhalteweg auswirken, und vergleichst deine
Ergebnisse mit Faustregeln, die du vielleicht schon aus der Fahrschule
oder dem Verkehrsunterricht kennst.
\end{beispielbox}

\begin{theoriebox}[Anwendung: Der Anhalteweg]
Eine der wichtigsten Alltagsanwendungen von Bewegungsgleichung 3: Wie
weit fährt ein Auto ab dem Moment, in dem ein Hindernis sichtbar wird,
bis es nach einer Vollbremsung steht? Dieser \textbf{\emph{Anhalteweg}}
$s_A$ setzt sich aus zwei Teilen zusammen:

\begin{itemize}[leftmargin=1.2cm,itemsep=0.25em]
  \item \textbf{\emph{Reaktionsweg}} $s_R$: die Strecke, die das Auto
  während der \textbf{\emph{Reaktionszeit}} $t_R$ zurücklegt, bevor die
  Fahrerin oder der Fahrer überhaupt zu bremsen beginnt. In dieser Zeit
  bleibt $v$ noch konstant bei $v_0$. Für eine nüchterne, aufmerksame
  Fahrerin oder einen nüchternen, aufmerksamen Fahrer wird meist
  $t_R\approx\SI{1}{\second}$ angenommen (in diesem Arbeitsblatt immer
  so, ausser wenn ausdrücklich etwas anderes angegeben ist).
  \item \textbf{\emph{Bremsweg}} $s_B$: die Strecke, die das Auto
  während der eigentlichen Bremsung zurücklegt, bis $v=0$ ist.
\end{itemize}

\begin{center}
\begin{tikzpicture}[xscale=1.15,yscale=0.32]
  \draw[->] (-0.3,0) -- (5.7,0) node[right] {\scriptsize $t$ in s};
  \draw[->] (0,-1) -- (0,23) node[above] {\scriptsize $v$ in $\si{\meter\per\second}$};
  \filldraw[titleblue!20] (0,0) rectangle (1,20);
  \filldraw[brakered!20] (1,0) -- (1,20) -- (5,0) -- cycle;
  \draw[thick,titleblue] (0,20) -- (1,20);
  \draw[thick,brakered] (1,20) -- (5,0);
  \draw[dashed,gray] (1,0) -- (1,20);
  \node at (0.5,9) {\scriptsize $s_R$};
  \node at (2.5,6) {\scriptsize $s_B$};
  \draw (1,0.5) -- (1,-0.5) node[below] {\scriptsize $t_R$};
  \node[left] at (0,20) {\scriptsize $v_0$};
\end{tikzpicture}
\end{center}

Wie schon beim $v$-$t$-Diagramm bekannt: Beide Teilwege entsprechen
Flächen unter der Kurve.
\begin{itemize}[leftmargin=1.2cm,itemsep=0.25em]
  \item $s_R$ ist ein \textbf{\emph{Rechteck}} (konstante Geschwindigkeit
  während $t_R$): $s_R = v_0\cdot t_R$.
  \item $s_B$ ist ein \textbf{\emph{Dreieck}}. Die Bremsverzögerung
  schreiben wir als \textbf{\emph{positive}} Zahl $a_B$ (den Betrag der
  eigentlich negativen Beschleunigung $a=-a_B$ beim Bremsen). Aus
  Bewegungsgleichung 3, mit Endgeschwindigkeit $v=0$:
  \[
    0 = v_0^2 + 2\cdot(-a_B)\cdot s_B
  \]
  Schritt 1: das Minuszeichen ausmultiplizieren:
  \[
    0 = v_0^2 - 2a_B\,s_B.
  \]
  Schritt 2: $2a_B\,s_B$ auf beiden Seiten addieren:
  \[
    2a_B\,s_B = v_0^2.
  \]
  Schritt 3: auf beiden Seiten durch $2a_B$ teilen:
  \[
    s_B = \frac{v_0^2}{2a_B}.
  \]
\end{itemize}

Addiert ergibt das den \textbf{\emph{Anhalteweg}}:
\[
  \boxed{s_A = s_R + s_B = v_0\,t_R + \frac{v_0^2}{2a_B}}
\]

\textbf{Beispiel:} $v_0=\SI{20}{\meter\per\second}$
($\approx\SI{72}{\kilo\meter\per\hour}$), $t_R=\SI{1}{\second}$,
$a_B=\SI{5}{\meter\per\second\squared}$ (nasse Strasse):
\[
  s_R = 20\cdot1=\SI{20}{\meter}, \qquad
  s_B = \frac{20^2}{2\cdot5} = \frac{400}{10}=\SI{40}{\meter},
\]
\[
  s_A = \SI{20}{\meter}+\SI{40}{\meter} = \SI{60}{\meter}.
\]
\end{theoriebox}

\begin{theoriebox}[Wird die Bremse rechtzeitig fertig? Das s-v-Diagramm]
Die Formel $s_A(v_0)=v_0\,t_R+\dfrac{v_0^2}{2a_B}$ zeigt: Der
Reaktionsweg wächst \textbf{\emph{linear}} mit $v_0$, der Bremsweg aber
\textbf{\emph{quadratisch}}. Um das sichtbar zu machen, tragen wir jetzt
genau wie in der Formel die Anfangsgeschwindigkeit $v_0$ waagrecht auf
(die unabhängige Grösse) und den zurückgelegten \textbf{\emph{Ort}} $s$
senkrecht (die abhängige Grösse, also $s$ als Funktion von $v_0$; mit
$t_R=\SI{1}{\second}$, $a_B=\SI{5}{\meter\per\second\squared}$, wie im
Beispiel oben):

\begin{center}
\begin{tikzpicture}[xscale=0.28,yscale=0.075]
  \draw[->] (-1,0) -- (33,0) node[right] {\scriptsize $v_0$ in $\si{\meter\per\second}$};
  \draw[->] (0,-4) -- (0,128) node[above] {\scriptsize $s$ in m};
  \fill[titleblue!18] (0,0) -- (30,0) -- (30,30) -- cycle;
  \fill[brakered!25,domain=0:30,samples=60]
    plot (\x,{\x+\x*\x/10}) -- (30,30) -- cycle;
  \draw[dotted,titleblue,thick] (0,0) -- (30,30);
  \draw[thick,brakered,domain=0:30,samples=60]
    plot (\x,{\x+\x*\x/10});
  \draw[thick,accentorange] (-1,60) -- (33,60);
  \node[accentorange,above] at (8,60) {\scriptsize Hindernis bei $s=\SI{60}{\meter}$};
  \draw[dashed,gray] (20,0) -- (20,60);
  \draw[dashed,gray] (0,60) -- (20,60);
  \filldraw[black] (20,60) circle (6pt);
  \node[titleblue] at (13,11) {\scriptsize $s_R=v_0\,t_R$};
  \node[brakered] at (24,80) {\scriptsize $s_B=\dfrac{v_0^2}{2a_B}$};
  \foreach \vv in {10,20,30} \draw (\vv,2) -- (\vv,-2) node[below] {\scriptsize \vv};
  \foreach \sv in {20,40,60,80,100,120} \draw (0.6,\sv) -- (-0.6,\sv) node[left] {\scriptsize \sv};
\end{tikzpicture}
\end{center}

Blau: der reine Reaktionsweg $s_R(v_0)=v_0\,t_R$, eine Gerade (die
gepunktete Linie). Rot: der zusätzliche Bremsweg, der zum Reaktionsweg
dazukommt, bis der volle Anhalteweg $s_A(v_0)$ erreicht ist (die obere
Kurve, die sich zunehmend nach oben krümmt, weil der Bremsweg
quadratisch wächst).

Angenommen, bei $s=\SI{60}{\meter}$ steht ein Hindernis (die orange
Linie). Für jede Anfangsgeschwindigkeit $v_0$, deren Kurve noch
unterhalb der orangen Linie bleibt, reicht der Anhalteweg aus: Das Auto
steht rechtzeitig. Beim markierten Punkt, $v_0=\SI{20}{\meter\per\second}$,
berührt die Kurve die orange Linie gerade noch: Das ist die
\textbf{\emph{höchste Geschwindigkeit}}, bei der das Auto genau noch
rechtzeitig steht (bestätigt das Beispiel der vorigen Theoriebox). Für
jedes $v_0>\SI{20}{\meter\per\second}$ liegt die Kurve oberhalb der
orangen Linie: Der Anhalteweg reicht nicht mehr aus, das Auto
kollidiert mit dem Hindernis, bevor es zum Stehen kommt.
\end{theoriebox}

\begin{questions}
\begin{taskbox}[Übung: Reaktionszeit unter Alkoholeinfluss]
\question[9] Alkohol verlängert die Reaktionszeit deutlich. Als
Beispielwert nehmen wir für eine leicht durch Alkohol beeinträchtigte
Fahrerin oder einen leicht beeinträchtigten Fahrer
$t_R=\SI{2}{\second}$, also doppelt so lang wie bei einer nüchternen
Person ($t_R=\SI{1}{\second}$). Untergrund und Hindernis bleiben
gleich: nasse Strasse ($a_B=\SI{5}{\meter\per\second\squared}$),
Hindernis bei $s=\SI{60}{\meter}$.

\begin{parts}
  \part[4] Berechne $s_R$, $s_B$ und $s_A$ für $v_0=\SI{20}{\meter\per\second}$
  mit $t_R=\SI{2}{\second}$.
  \Antwortraum{4}
  \begin{solution}
  $s_R = 20\cdot2=\SI{40}{\meter}$.\\
  $s_B = \dfrac{20^2}{2\cdot5}=\SI{40}{\meter}$ (unverändert, $a_B$ ist
  gleich geblieben).\\
  $s_A = \SI{40}{\meter}+\SI{40}{\meter}=\SI{80}{\meter}$.
  \end{solution}

  \part[2] Bei $v_0=\SI{20}{\meter\per\second}$ stand die nüchterne
  Fahrerin aus der Theoriebox gerade noch rechtzeitig (Beispiel oben).
  Reicht das auch für die beeinträchtigte Person? Begründe mit deinem
  Ergebnis aus a).
  \Antwortraum{2}
  \begin{solution}
  Nein: $s_A=\SI{80}{\meter}$ ist grösser als die
  $\SI{60}{\meter}$ bis zum Hindernis. Bei genau derselben Geschwindigkeit,
  die für die nüchterne Person gerade noch sicher war, kollidiert die
  beeinträchtigte Person, weil allein die längere Reaktionszeit
  $\SI{20}{\meter}$ mehr Anhalteweg braucht.
  \end{solution}

  \part[3] Zeichne unten das $s$-$v$-Diagramm für die beeinträchtigte
  Person (mit $t_R=\SI{2}{\second}$), rechts direkt neben dem bereits
  bekannten Diagramm für die nüchterne Person ($t_R=\SI{1}{\second}$).
  Markiere auch hier den Punkt, an dem die Kurve die orange
  Hindernislinie erreicht.

  \begin{center}
  \begin{minipage}{0.47\linewidth}\centering
  {\scriptsize \textbf{Nüchtern, $t_R=\SI{1}{\second}$} (zur Erinnerung)}\\[0.3em]
  \begin{tikzpicture}[xscale=0.13,yscale=0.033]
    \draw[->] (-1,0) -- (33,0) node[right] {\scriptsize $v_0$};
    \draw[->] (0,-4) -- (0,158) node[above] {\scriptsize $s$};
    \draw[dotted,titleblue,thick] (0,0) -- (30,30);
    \draw[thick,brakered,domain=0:30,samples=60] plot (\x,{\x+\x*\x/10});
    \draw[thick,accentorange] (-1,60) -- (33,60);
    \filldraw[black] (20,60) circle (5pt);
    \foreach \vv in {10,20,30} \draw (\vv,2) -- (\vv,-2) node[below] {\scriptsize \vv};
    \foreach \sv in {60,120} \draw (0.6,\sv) -- (-0.6,\sv) node[left] {\scriptsize \sv};
  \end{tikzpicture}
  \end{minipage}\hfill
  \begin{minipage}{0.47\linewidth}\centering
  {\scriptsize \textbf{Beeinträchtigt, $t_R=\SI{2}{\second}$} (deine Zeichnung)}\\[0.3em]
  \begin{tikzpicture}[xscale=0.13,yscale=0.033]
    \draw[->] (-1,0) -- (33,0) node[right] {\scriptsize $v_0$};
    \draw[->] (0,-4) -- (0,158) node[above] {\scriptsize $s$};
    \draw[thick,accentorange] (-1,60) -- (33,60);
    \foreach \vv in {10,20,30} \draw (\vv,2) -- (\vv,-2) node[below] {\scriptsize \vv};
    \foreach \sv in {60,120} \draw (0.6,\sv) -- (-0.6,\sv) node[left] {\scriptsize \sv};
  \end{tikzpicture}
  \end{minipage}
  \end{center}
  \begin{solution}
  \begin{center}
  \begin{tikzpicture}[xscale=0.13,yscale=0.033]
    \draw[->] (-1,0) -- (33,0) node[right] {\scriptsize $v_0$ in $\si{\meter\per\second}$};
    \draw[->] (0,-4) -- (0,158) node[above] {\scriptsize $s$ in m};
    \fill[titleblue!18] (0,0) -- (30,0) -- (30,60) -- cycle;
    \fill[brakered!25,domain=0:30,samples=60] plot (\x,{2*\x+\x*\x/10}) -- (30,60) -- cycle;
    \draw[dotted,titleblue,thick] (0,0) -- (30,60);
    \draw[thick,brakered,domain=0:30,samples=60] plot (\x,{2*\x+\x*\x/10});
    \draw[thick,accentorange] (-1,60) -- (33,60);
    \draw[dashed,gray] (16.46,0) -- (16.46,60);
    \filldraw[black] (16.46,60) circle (5pt);
    \foreach \vv in {10,20,30} \draw (\vv,2) -- (\vv,-2) node[below] {\scriptsize \vv};
    \foreach \sv in {60,120} \draw (0.6,\sv) -- (-0.6,\sv) node[left] {\scriptsize \sv};
  \end{tikzpicture}
  \end{center}
  Die Kurve steigt jetzt schon zu Beginn steiler an (grössere
  Anfangssteigung $=t_R=2$) und erreicht die orange Linie schon bei
  $v_0\approx\SI{16.5}{\meter\per\second}$, also bei einer deutlich
  \emph{tieferen} Geschwindigkeit als bei der nüchternen Person
  ($\SI{20}{\meter\per\second}$). Die \glqq sichere\grqq{}
  Höchstgeschwindigkeit sinkt also allein durch die längere
  Reaktionszeit.
  \end{solution}
\end{parts}
\end{taskbox}
\end{questions}

\begin{questions}
\begin{taskbox}[Übung: Bremsweg und Untergrund]
\question[10] Die Bremsverzögerung $a_B$ hängt stark vom Untergrund ab.
Typische Werte:
\begin{center}
{\small
\begin{tabular}{lc}
\toprule
\textbf{Untergrund} & \textbf{$a_B$ in $\si{\meter\per\second\squared}$} \\
\midrule
trockene Strasse & 8 \\
nasse Strasse & 5 \\
Neuschnee & 2.5 \\
Glatteis & 1.5 \\
\bottomrule
\end{tabular}
}
\end{center}

\begin{parts}
  \part[4] Trägt man den Bremsweg $s_B(v_0)=\dfrac{v_0^2}{2a_B}$ für
  jeden dieser vier Werte über $v_0$ auf, ergeben sich vier verschiedene
  Parabeln. Berechne für jeden Untergrund $s_B$ bei
  $v_0=\SI{10}{\meter\per\second}$ und bei $v_0=\SI{20}{\meter\per\second}$
  (das reicht für eine grobe Skizze) und zeichne alle vier Parabeln
  näherungsweise ins Diagramm.

  \begin{center}
  \begin{tikzpicture}[xscale=0.5,yscale=0.055]
    \draw[->] (-0.6,0) -- (21,0) node[right] {\scriptsize $v_0$ in $\si{\meter\per\second}$};
    \draw[->] (0,-4) -- (0,145) node[above] {\scriptsize $s_B$ in m};
    \foreach \xv in {5,10,15,20} \draw (\xv,3) -- (\xv,-3) node[below] {\scriptsize \xv};
    \foreach \yv in {25,50,75,100,125} \draw (0.4,\yv) -- (-0.4,\yv) node[left] {\scriptsize \yv};
  \end{tikzpicture}
  \end{center}
  \begin{solution}
  $v_0=\SI{10}{\meter\per\second}$: trocken $\SI{6.25}{\meter}$, nass
  $\SI{10}{\meter}$, Neuschnee $\SI{20}{\meter}$, Glatteis
  $\SI{33.3}{\meter}$.\\
  $v_0=\SI{20}{\meter\per\second}$: trocken $\SI{25}{\meter}$, nass
  $\SI{40}{\meter}$, Neuschnee $\SI{80}{\meter}$, Glatteis
  $\SI{133.3}{\meter}$.
  \begin{center}
  \begin{tikzpicture}[xscale=0.5,yscale=0.055]
    \draw[->] (-0.6,0) -- (21,0) node[right] {\scriptsize $v_0$ in $\si{\meter\per\second}$};
    \draw[->] (0,-4) -- (0,145) node[above] {\scriptsize $s_B$ in m};
    \draw[thick,titleblue] plot[domain=0:20,samples=50] (\x,{\x*\x/16})
      node[right] {\scriptsize trocken};
    \draw[thick,brakered] plot[domain=0:20,samples=50] (\x,{\x*\x/10})
      node[right] {\scriptsize nass};
    \draw[thick,accentpurple] plot[domain=0:20,samples=50] (\x,{\x*\x/5})
      node[right] {\scriptsize Neuschnee};
    \draw[thick,groupteal] plot[domain=0:20,samples=50] (\x,{\x*\x/3})
      node[right] {\scriptsize Glatteis};
    \foreach \xv in {5,10,15,20} \draw (\xv,3) -- (\xv,-3) node[below] {\scriptsize \xv};
    \foreach \yv in {25,50,75,100,125} \draw (0.4,\yv) -- (-0.4,\yv) node[left] {\scriptsize \yv};
  \end{tikzpicture}
  \end{center}
  Bei $v_0=\SI{20}{\meter\per\second}$ liegt der Bremsweg auf Glatteis
  mehr als fünfmal so hoch wie auf trockener Strasse, obwohl die
  Geschwindigkeit exakt gleich ist: Nur die Bremsverzögerung $a_B$, also
  die Steilheit der jeweiligen Parabel, hat sich geändert.
  \end{solution}

  \part[3] Bei Aquaplaning bildet sich zwischen Reifen und Fahrbahn ein
  Wasserfilm, sodass die Reifen den Kontakt zur Strasse verlieren.
  Erkläre mithilfe der Formel für $s_B$, weshalb das besonders
  gefährlich ist.
  \Antwortraum{3}
  \begin{solution}
  Ohne Reibungskontakt zwischen Reifen und Strasse kann das Auto kaum
  noch gebremst (und auch kaum noch gelenkt) werden: $a_B$ sinkt dabei
  stark, im Extremfall fast auf $\SI{0}{\meter\per\second\squared}$.
  Weil $s_B=\dfrac{v_0^2}{2a_B}$ ist, wächst der Bremsweg, wenn $a_B$
  gegen null geht, über alle Grenzen: Im schlimmsten Fall bremst das
  Auto praktisch gar nicht mehr ab.
  \end{solution}

  \part[3] Erkläre, weshalb Glatteis auch im Vergleich zu Neuschnee
  besonders gefährlich ist (nutze die Tabelle und die Formel für $s_B$).
  \Antwortraum{3}
  \begin{solution}
  Glatteis hat mit $a_B=\SI{1.5}{\meter\per\second\squared}$ die
  kleinste Bremsverzögerung aller vier Untergründe, kleiner sogar als
  Neuschnee ($\SI{2.5}{\meter\per\second\squared}$). Weil $s_B$
  umgekehrt proportional zu $a_B$ ist, wächst der Bremsweg bei
  kleinerem $a_B$ überproportional an; gleichzeitig wächst $s_B$ auch
  bei gleichbleibendem $a_B$ quadratisch mit $v_0$. Beide Effekte
  verstärken sich gegenseitig, weshalb schon eine mässige
  Geschwindigkeit auf Glatteis zu einem um ein Vielfaches längeren
  Bremsweg führen kann als auf trockener Strasse.
  \end{solution}
\end{parts}
\end{taskbox}
\end{questions}

\begin{hinweisbox}[Rückwärts rechnen: $v_0$ aus dem Bremsweg]
Manchmal ist nicht $v_0$ gegeben, sondern der Bremsweg, und $v_0$ ist
gesucht, zum Beispiel: \glqq Wie schnell war ein Auto höchstens
unterwegs, wenn seine Bremsspur $\SI{40}{\meter}$ lang ist?\grqq{} Das
bedeutet, $v_0^2=2a_B\,s_B$ nach $v_0$ aufzulösen. Das ist
\textbf{\emph{einfacher}} als die Mitternachtsformel: Es gibt hier
keinen linearen Term (keinen $B$-Term wie bei $At^2+Bt+C$), also reicht
es, auf beiden Seiten die \textbf{\emph{Wurzel}} zu ziehen:
\[
  v_0^2 = 2a_B\,s_B
  \quad\Longrightarrow\quad
  v_0 = \sqrt{2a_B\,s_B}.
\]
\textbf{Beispiel:} $a_B=\SI{5}{\meter\per\second\squared}$,
$s_B=\SI{40}{\meter}$:
\[
  v_0 = \sqrt{2\cdot5\cdot40} = \sqrt{400} = \SI{20}{\meter\per\second}.
\]
Achtung: Das funktioniert nur, wenn im Term wirklich nur $v_0^2$
vorkommt, kein einzelnes $v_0$ daneben (wie es beim vollen Anhalteweg
$s_A(v_0)=v_0 t_R+\frac{v_0^2}{2a_B}$ der Fall wäre). Dann bleibt es bei
der Mitternachtsformel, siehe unten.
\end{hinweisbox}

\begin{questions}
\begin{taskbox}[Übung: Anhalteweg berechnen]
\question[14] Ein Auto fährt mit $v_0=\SI{25}{\meter\per\second}$
($\approx\SI{90}{\kilo\meter\per\hour}$) auf einer nassen Strasse
($a_B=\SI{5}{\meter\per\second\squared}$). Die Reaktionszeit der
Fahrerin beträgt $t_R=\SI{1}{\second}$.

\begin{parts}
  \part[2] Berechne den Reaktionsweg $s_R$.
  \Antwortraum{2}
  \begin{solution}
  $s_R = v_0\cdot t_R = 25\cdot1=\SI{25}{\meter}$.
  \end{solution}

  \part[3] Berechne den Bremsweg $s_B$ mit Bewegungsgleichung 3
  ($s_B=\dfrac{v_0^2}{2a_B}$).
  \Antwortraum{3}
  \begin{solution}
  $s_B = \dfrac{25^2}{2\cdot5} = \dfrac{625}{10} = \SI{62.5}{\meter}$.
  \end{solution}

  \part[2] Wie gross ist der gesamte Anhalteweg $s_A$?
  \Antwortraum{2}
  \begin{solution}
  $s_A = s_R+s_B = \SI{25}{\meter}+\SI{62.5}{\meter}=\SI{87.5}{\meter}$.
  \end{solution}

  \part[3] Auf trockener Strasse wäre $a_B=\SI{8}{\meter\per\second\squared}$.
  Berechne den Anhalteweg für diesen Fall (gleiches $v_0$, gleiches
  $t_R$) und vergleiche mit Teilaufgabe c).
  \Antwortraum{3}
  \begin{solution}
  $s_B = \dfrac{625}{16}\approx\SI{39.1}{\meter}$.
  $s_A = \SI{25}{\meter}+\SI{39.1}{\meter}\approx\SI{64.1}{\meter}$.
  Das ist rund $\SI{23}{\meter}$ kürzer als auf nasser Strasse, obwohl
  $v_0$ und $t_R$ gleich geblieben sind: Nur der Bremsweg hängt vom
  Untergrund ab, der Reaktionsweg nicht.
  \end{solution}

  \part[4] Auf einer anderen Fahrt hinterlässt ein Auto auf trockener
  Strasse ($a_B=\SI{8}{\meter\per\second\squared}$) eine Bremsspur von
  $s_B=\SI{72}{\meter}$. Wie schnell war das Auto zu Beginn der Bremsung
  unterwegs? Nutze die Methode aus dem Kasten oben.
  \Antwortraum{4}
  \begin{solution}
  \[
    v_0 = \sqrt{2a_B\,s_B} = \sqrt{2\cdot8\cdot72} = \sqrt{1152}
    \approx \SI{33.9}{\meter\per\second} \approx \SI{122}{\kilo\meter\per\hour}.
  \]
  \end{solution}
\end{parts}
\end{taskbox}
\end{questions}

\begin{theoriebox}[Faustregeln aus der Fahrschule]
Neben den exakten Bewegungsgleichungen gibt es aus der Fahrschule
bekannte Faustregeln, die sich im Kopf rechnen lassen. Zwei davon
betreffen den \textbf{\emph{Sicherheitsabstand}} zu einem
\emph{vorausfahrenden} Fahrzeug, also einen anderen Fall als unser
Hindernis, das stillsteht:
\begin{itemize}[leftmargin=1.2cm,itemsep=0.25em]
  \item \textbf{\emph{2-Sekunden-Regel}}: Der Abstand zum
  vorausfahrenden Fahrzeug soll so gross sein, dass mindestens
  $\SI{2}{\second}$ vergehen, bis man selbst denselben Punkt erreicht:
  $s_{2\text{s}} = v_0\cdot\SI{2}{\second}$.
  \item \textbf{\emph{Halber-Tacho-Regel}}: Sicherheitsabstand in
  Metern $\approx$ die Hälfte der Tachoanzeige in
  $\si{\kilo\meter\per\hour}$: $s_{\text{Tacho}}
  \approx \dfrac{v_0\,[\si{\kilo\meter\per\hour}]}{2}$.
\end{itemize}
Eine ältere, ebenfalls im Kopf rechenbare Faustformel betrifft dagegen
speziell den \textbf{\emph{Bremsweg}} allein, nicht den ganzen
Anhalteweg:
\[
  s_B\,[\text{m}] \approx \left(\frac{v_0\,[\si{\kilo\meter\per\hour}]}{10}\right)^{\!2}.
\]
\textbf{Wichtig:} Diese drei Faustregeln sind praktische Abschätzungen
für den Alltag, kein Ersatz für die exakten Bewegungsgleichungen. Die
ersten beiden schätzen den Abstand zu einem \emph{fahrenden} Fahrzeug
ab, nicht den vollen Anhalteweg vor einem \emph{stillstehenden}
Hindernis. Die dritte deckt nur den Bremsweg ab, nicht den
Reaktionsweg.
\end{theoriebox}

\begin{questions}
\begin{taskbox}[Übung: Faustregeln im Vergleich]
\question[9] Betrachte wieder das Auto von vorhin:
$v_0=\SI{25}{\meter\per\second}$ ($\approx\SI{90}{\kilo\meter\per\hour}$),
nasse Strasse.

\begin{parts}
  \part[3] Berechne alle drei Faustregel-Werte für dieses Auto:
  $s_{2\text{s}}$, $s_{\text{Tacho}}$ und $s_B$ nach der Faustformel.
  \Antwortraum{3}
  \begin{solution}
  $s_{2\text{s}} = 25\cdot2=\SI{50}{\meter}$.\\
  $s_{\text{Tacho}} \approx 90/2=\SI{45}{\meter}$.\\
  $s_B\text{ (Faustformel)} \approx (90/10)^2=\SI{81}{\meter}$.
  \end{solution}

  \part[3] Vergleiche $s_B$ nach der Faustformel mit dem exakten
  Bremsweg aus der vorigen Übung ($s_B=\SI{62.5}{\meter}$ auf nasser
  Strasse). Was passiert, wenn jemand die Faustformel fälschlicherweise
  für den \emph{gesamten} Anhalteweg hält statt nur für den Bremsweg?
  Wie gross wäre der Fehler?
  \Antwortraum{3}
  \begin{solution}
  Die Faustformel ($\SI{81}{\meter}$) liegt sogar über dem exakten
  Bremsweg auf nasser Strasse ($\SI{62.5}{\meter}$): Als reine
  Bremsweg-Abschätzung ist sie eher grosszügig. Hält man die
  $\SI{81}{\meter}$ aber fälschlich für den \emph{gesamten} Anhalteweg,
  fehlt der Reaktionsweg von $\SI{25}{\meter}$ komplett: Der
  tatsächliche Anhalteweg ($\SI{87.5}{\meter}$) ist grösser als die
  $\SI{81}{\meter}$ der Faustformel. Genau diese Verwechslung von
  Bremsweg und gesamtem Anhalteweg ist ein klassischer und gefährlicher
  Denkfehler.
  \end{solution}

  \part[3] Angenommen, das Auto folgt einem anderen Fahrzeug mit
  $\SI{45}{\meter}$ Abstand (Halber-Tacho-Regel eingehalten). Reicht
  dieser Abstand aus, um vor einem plötzlich \emph{stillstehenden}
  Hindernis (zum Beispiel ein liegen gebliebenes Auto) rechtzeitig zu
  stehen? Begründe mit dem exakten Anhalteweg $s_A=\SI{87.5}{\meter}$.
  \Antwortraum{3}
  \begin{solution}
  Nein: $\SI{45}{\meter}$ sind deutlich weniger als der exakte
  Anhalteweg von $\SI{87.5}{\meter}$. Die Halber-Tacho-Regel (wie auch
  die 2-Sekunden-Regel) ist für den Abstand zu einem \emph{fahrenden}
  Fahrzeug gedacht, das im Ernstfall selbst auch bremst; sie ersetzt
  nicht den vollen Anhalteweg vor einem stillstehenden Hindernis. Wer
  diese beiden Situationen verwechselt, unterschätzt den nötigen
  Abstand erheblich.
  \end{solution}
\end{parts}
\end{taskbox}
\end{questions}

\begin{theoriebox}[Quadratische Gleichungen lösen: die Mitternachtsformel]
Löst man die Ortsgleichung $s(t)=s_0+v_0(t-t_0)+\frac{1}{2}a(t-t_0)^2$
nach $t$ auf (statt nach $s$), steht $t$ im Quadrat: Du bekommst eine
\textbf{\emph{quadratische Gleichung}}. Für eine allgemeine quadratische
Gleichung in der Form
\[
  A\,t^2 + B\,t + C = 0
\]
(dieselben Buchstaben $A,B,C$ wie bei der allgemeinen Parabelgleichung
$f(t)=At^2+Bt+C$ aus Teil 1, jetzt aber als Gleichung, nicht als
Funktion) liefert die
\textbf{\emph{Mitternachtsformel}} beide Lösungen auf einen Schlag:
\[
  \boxed{t = \frac{-B \pm \sqrt{B^2-4AC}}{2A}}
\]
Das $\pm$ bedeutet: Du rechnest die Formel \textbf{\emph{zweimal}}, einmal
mit $+$ und einmal mit $-$ vor der Wurzel, und erhältst dabei bis zu zwei
verschiedene Lösungen für $t$.

\textbf{Beispiel (rein mathematisch, ohne Physik):} Löse
$2t^2+4t-6=0$. Hier ist $A=2$, $B=4$, $C=-6$:
\[
  t = \frac{-4\pm\sqrt{4^2-4\cdot2\cdot(-6)}}{2\cdot2}
  = \frac{-4\pm\sqrt{64}}{4} = \frac{-4\pm8}{4}.
\]
Das ergibt $t_1=1$ und $t_2=-3$.

\textbf{In der Physik:} Um die Ortsgleichung auf diese Form zu bringen,
bringst du einfach alles auf eine Seite: Aus
$s_0+v_0(t-t_0)+\frac{1}{2}a(t-t_0)^2 = s_{\text{gesucht}}$ wird (mit
$t_0=0$)
\[
  \underbrace{\tfrac{1}{2}a}_{A}\,t^2 + \underbrace{v_0}_{B}\,t + \underbrace{(s_0-s_{\text{gesucht}})}_{C} = 0,
\]
und du kannst $A$, $B$, $C$ direkt in die Mitternachtsformel einsetzen.
Du brauchst dieses Verfahren immer dann, wenn eine Aufgabe nach der Zeit
$t$ fragt, aber die Ortsgleichung (nicht die Geschwindigkeitsgleichung)
die passende Formel ist, siehe Checkliste in Teil 1.

Eine quadratische Gleichung hat mathematisch bis zu \textbf{\emph{zwei}}
Lösungen (wegen des $\pm$). In der Physik ist aber meistens nur eine
davon \textbf{\emph{physikalisch sinnvoll}}.

\textbf{Wichtig:} Verwirf nicht automatisch nur negative Zeiten. Prüfe
bei \emph{jeder} Lösung, ob sie in der Situation überhaupt Sinn ergibt.
Eine Lösung kann auch aus einem anderen Grund unphysikalisch sein, zum
Beispiel weil die Bewegung zu diesem Zeitpunkt laut Aufgabenstellung
längst vorbei ist (siehe Beispiel unten).
\end{theoriebox}

\begin{questions}
\begin{taskbox}[Zusatzaufgabe: Die kritische Geschwindigkeit]
\question[6] Auf einer nassen Strasse ($a_B=\SI{5}{\meter\per\second\squared}$,
$t_R=\SI{1}{\second}$) steht bei $s=\SI{60}{\meter}$ ein Hindernis, wie
im $s$-$v$-Diagramm weiter oben. Diesmal ist $v_0$ die gesuchte Grösse.

\begin{parts}
  \part[2] Stelle die Gleichung $s_A(v_0)=v_0\,t_R+\dfrac{v_0^2}{2a_B}=\SI{60}{\meter}$
  mit den gegebenen Zahlenwerten auf und bring sie in die Form
  $A\,v_0^2+B\,v_0+C=0$.
  \Antwortraum{2}
  \begin{solution}
  \[
    v_0 + \frac{v_0^2}{10} = 60
  \]
  Schritt 1: mit dem gemeinsamen Nenner $10$ durchmultiplizieren, um den
  Bruch zu beseitigen:
  \[
    10v_0 + v_0^2 = 600.
  \]
  Schritt 2: alles auf eine Seite bringen (Standardform $A\,v_0^2+B\,v_0+C=0$):
  \[
    v_0^2 + 10v_0 - 600 = 0.
  \]
  Also $A=1$, $B=10$, $C=-600$. (Hier steht sowohl ein $v_0^2$- als auch
  ein einzelner $v_0$-Term, deshalb reicht die einfache Wurzel aus dem
  Kasten oben \emph{nicht}: Das ist ein Fall für die Mitternachtsformel.)
  \end{solution}

  \part[4] Löse mit der Mitternachtsformel nach $v_0$ auf und verwirf die
  unphysikalische Lösung.
  \Antwortraum{4}
  \begin{solution}
  \[
    v_0 = \frac{-10\pm\sqrt{10^2-4\cdot1\cdot(-600)}}{2\cdot1}.
  \]
  \[
    v_0 = \frac{-10\pm\sqrt{100+2400}}{2} = \frac{-10\pm\sqrt{2500}}{2}.
  \]
  \[
    v_0 = \frac{-10\pm50}{2}.
  \]
  Das ergibt $v_0=\SI{20}{\meter\per\second}$ oder $v_0=-\SI{30}{\meter\per\second}$.
  Eine negative Anfangsgeschwindigkeit ergibt hier keinen Sinn, also
  bleibt $v_0=\SI{20}{\meter\per\second}$: genau der Punkt, den das
  $s$-$v$-Diagramm oben markiert.
  \end{solution}
\end{parts}
\end{taskbox}
\end{questions}

\begin{questions}
\begin{taskbox}[Übung: E-Trottinett bremst, zweiter Teil]
\question[6] Dasselbe E-Trottinett wie im Beispiel in Teil 1: $v_0=\SI{6}{\meter\per\second}$,
$a=-\SI{2}{\meter\per\second\squared}$, positive Richtung = Fahrtrichtung.

\begin{parts}
  \part[2] Nach welcher Zeit $t_{\text{stopp}}$ steht das Trottinett
  (also $v=0$)? Prüfe dein Ergebnis anschliessend, indem du die
  Bremsstrecke aus dem Beispiel in Teil 1 ($\Delta s=\SI{9}{\meter}$) über
  $\Delta s = v_0\cdot t_{\text{stopp}} + \frac{1}{2}a\,t_{\text{stopp}}^2$
  nachrechnest.
  \Antwortraum{2}
  \begin{solution}
  $0 = 6 - 2\cdot t_{\text{stopp}} \Rightarrow t_{\text{stopp}}=\SI{3}{\second}$.
  Probe: $\Delta s = 6\cdot3 - \frac{1}{2}\cdot2\cdot3^2 = 18-9=\SI{9}{\meter}$,
  stimmt mit dem Beispiel in Teil 1 überein.
  Ab diesem Zeitpunkt gilt die Formel mit konstantem $a$ nicht mehr, das
  Trottinett bleibt einfach stehen.
  \end{solution}

  \part[4] Nach welcher Zeit $t$ hat das Trottinett eine Strecke von
  $\Delta s=\SI{5}{\meter}$ zurückgelegt? Löse mit der Ortsgleichung
  ($s_0=0$, $t_0=0$) und diskutiere \textbf{beide} Lösungen der
  quadratischen Gleichung.
  \Antwortraum{4}
  \begin{solution}
  Zuerst alles auf eine Seite bringen (Standardform $At^2+Bt+C=0$):
  \[
    5 = 6t - t^2 \quad\Longrightarrow\quad t^2-6t+5=0.
  \]
  Koeffizienten ablesen: $A=1$, $B=-6$, $C=5$. In die Mitternachtsformel
  einsetzen:
  \[
    t = \frac{-B\pm\sqrt{B^2-4AC}}{2A} = \frac{6\pm\sqrt{(-6)^2-4\cdot1\cdot5}}{2\cdot1}.
  \]
  \[
    t = \frac{6\pm\sqrt{36-20}}{2} = \frac{6\pm\sqrt{16}}{2}.
  \]
  \[
    t = \frac{6\pm4}{2}.
  \]
  Das ergibt $t_1=\frac{6+4}{2}=\SI{5}{\second}$ und
  $t_2=\frac{6-4}{2}=\SI{1}{\second}$.
  Mathematisch gibt es also zwei Lösungen: $t=\SI{1}{\second}$ und
  $t=\SI{5}{\second}$. Aus Teilaufgabe a) weisst du aber, dass das
  Trottinett bereits bei $t_{\text{stopp}}=\SI{3}{\second}$ steht: Die
  Formel mit konstantem $a$ gilt also nur für $0\leq t\leq\SI{3}{\second}$.
  $t=\SI{5}{\second}$ liegt \emph{ausserhalb} dieses gültigen Bereichs
  (nicht etwa negativ!) und ist deshalb zu verwerfen. Die physikalisch
  richtige Lösung ist $t=\SI{1}{\second}$.
  \end{solution}
\end{parts}
\end{taskbox}
\end{questions}

\end{document}
